\chapter*{Resumen}\addcontentsline{toc}{chapter}{Resumen}

Este borrador presenta el estado actual del proyecto ``Astro Denoise Benchmark'' y documenta el alcance hasta el estado del arte y los datos preliminares. Se describen el contexto cient\'ifico, la motivaci\'on inspiracional para una ciencia m\'as abierta y transparente, el planteamiento del problema y los datasets DC2 que se usar\'an para evaluaci\'on. Adem\'as, se documenta el procedimiento reproducible de selecci\'on bibliogr\'afica potenciado por aprendizaje activo y la integraci\'on inicial de referencias clave que definen el panorama actual del denoising astron\'omico.

El objetivo del documento es ofrecer una base t\'ecnicamente s\'olida y reproducible para comparar m\'etodos de denoising (enfoc\'andose en BM3D y U-Net) mediante m\'etricas orientadas a la ciencia \textemdash completitud y pureza\textemdash en im\'agenes simuladas del Vera Rubin Observatory (DC2). Este informe no contiene a\'un resultados experimentales finales: presenta el alcance, los datos disponibles y la metodolog\'ia propuesta para las pr\'oximas etapas del benchmark.

Se han versionado los scripts, los datos y las fichas bibliogr\'aficas para facilitar la auditor\'ia y la repetici\'on del estudio por parte de la comunidad cient\'ifica.

\bigskip

\textbf{Palabras clave:} denoising; im\'agenes astron\'omicas; DC2; BM3D; U-Net; completitud; pureza; Vera Rubin Observatory; reproducibilidad.

\bigskip

\textbf{Keywords:} denoising; astronomical images; DC2; BM3D; U-Net; completeness; purity; Vera Rubin Observatory; reproducibility.
